% Formation Setup

% \begin{center}
% \textbf{\large Optimal Formation Geometry}
% \end{center}

\begin{center}
    \begin{adjustbox}{width=0.8\columnwidth}
        \begin{tikzpicture}

            % two leader nodes
            \node (L0) at (-5,0) {\includegraphics[width=1.3cm]{figures/lolo_top.png}};
            \node[below=0cm of L0] {\(\mathbf{L_0}\)}; % Label for L_0

            \node (L1) at (5,0) {\includegraphics[width=1.3cm]{figures/lolo_top.png}};
            \node[below=0cm of L1] {\(\mathbf{L_1}\)}; % Label for L_1
                    
            % five followers symmetrically placed at 90° angles
            \node (f1) at ({-2.5 * sqrt(3)},-2.5) {\includegraphics[width=0.3cm]{figures/sam_top.png}};
            \node[below=0cm of f1] {\(\mathbf{f_0}\)}; % Label for f_1
            
            \node (f2) at (-2.5,{-2.5 * sqrt(3)}) {\includegraphics[width=0.3cm]{figures/sam_top.png}};
            \node[below=0cm of f2] {\(\mathbf{f_1}\)}; % Label for f_2
            
            \node (f3) at (0,-5) {\includegraphics[width=0.3cm]{figures/sam_top.png}};
            \node[below=0cm of f3] {\(\mathbf{f_2}\)}; % Label for f_3
            
            \node (f4) at (2.5,{-2.5 * sqrt(3)}) {\includegraphics[width=0.3cm]{figures/sam_top.png}};
            \node[below=0cm of f4] {\(\mathbf{f_3}\)}; % Label for f_4
            
            \node (f5) at ({2.5 * sqrt(3)},-2.5) {\includegraphics[width=0.3cm]{figures/sam_top.png}};
            \node[below=0cm of f5] {\(\mathbf{f_4}\)}; % Label for f_5

            % draw lines between leaders and followers
            \draw[dotted, thin] (L0.center) -- (f1.center);
            \draw[dotted, thin] (L1.center) -- (f1.center);
            \pic [draw, angle radius=5mm, angle eccentricity=1.5, "$90^\circ$", anchor=center, color=red] {right angle = L1--f1--L0}; 

            \draw[dotted, thin] (L0.center) -- (f2.center);
            \draw[dotted, thin] (L1.center) -- (f2.center);
            \pic [draw, angle radius=5mm, angle eccentricity=1.5, "$90^\circ$", anchor=center, color=red] {right angle = L1--f2--L0}; 

            \draw[dotted, thin] (L0.center) -- (f3.center);
            \draw[dotted, thin] (L1.center) -- (f3.center);
            \pic [draw, angle radius=5mm, angle eccentricity=1.5, "$90^\circ$", anchor=center, color=red] {right angle = L1--f3--L0}; 

            \draw[dotted, thin] (L0.center) -- (f4.center);
            \draw[dotted, thin] (L1.center) -- (f4.center);
            \pic [draw, angle radius=5mm, angle eccentricity=1.5, "$90^\circ$", anchor=center, color=red] {right angle = L1--f4--L0}; 

            \draw[dotted, thin] (L0.center) -- (f5.center);
            \draw[dotted, thin] (L1.center) -- (f5.center);
            \pic [draw, angle radius=5mm, angle eccentricity=1.5, "$90^\circ$", anchor=center, color=red] {right angle = L1--f5--L0}; 

            % semicircular arc with leaders on diameter (dashed, thick, orange)
            \draw[dashed, thick, orange] (L0.center) arc[start angle=180,end angle=360,radius=5cm];

        \end{tikzpicture}
    \end{adjustbox}
    \captionof{figure}{\footnotesize TUPER Formation Geometry}
\end{center}

\begin{tcolorbox}[colback=waspgrey!10, colframe=waspgrey, boxrule=0.5pt]
\textbf{Why Semi-Circle?}
\begin{itemize}[leftmargin=1em, itemsep=0.05em]
    \item Minimize GDOP \cite{taraldsen2011underwater}
    \item Balance b/w coverage and accuracy
    \item Scalable to arbitrary number of followers
\end{itemize}
\end{tcolorbox}
